\section{Two Grounding Models}
\label{sec:bridges}

Two concrete models demonstrate that our abstract theory is 
non-degenerate.  One shows every
classical measure \emph{contains} an intuitionistic probability.
The other shows the intuitionistic theory cannot collapse back
without solving the halting problem.

\subsection{Measures through the Interior Operator}
\label{sec:interior-bridge}

G\"odel--McKinsey--Tarski interpret intuitionistic logic in modal
$\mathsf{S4}$: propositions are \emph{open} sets, and
the modal necessity operator $\Box$ of $\mathsf{S4}$ is interpreted
as the topological interior, written $\mathrm{int}(\cdot)$
below~\cite{godel1933,mckinsey-tarski1944}.
Our valuations quantify that same translation:

\begin{theorem}[\texttt{Measure.toValuationOpens}]
Let $\mu$ be a probability measure on a space $X$, $A \subseteq X$, 
and define $v(A) = \mu(\mathrm{int}\,A)$.  Since
$\mathrm{int}\,U = U$ for open $U$, $v$ restricts on the frame
$\mathcal{O}(X)$ to $U \mapsto \mu(U)$, and this restriction is a
valuation.
\end{theorem}

\noindent
Modularity is $\mu(U \cup V) + \mu(U \cap V) = \mu(U) + \mu(V)$.
The slack of an open $U$ is then $1 - \mu(U) -
\mu(\mathrm{int}\,U^{\complement})$, which is the measure of the
boundary $\partial U$.  We can now prove that generally, slack is 
exactly the boundary mass:

\begin{theorem}\strut\\
{\upshape(\texttt{toValuationOpens\_slack\_eq\_frontier})}\\
For any probability measure $\mu$ on a topological space $X$ and any
open $U$, $\slack(U) = \mu(\partial U)$.
\end{theorem}

\noindent
The Heyting complement of $U$ in $\mathcal{O}(X)$ is the interior of
its set complement. So, the undecided region, the complement of
$U \sqcup \neg U$, is exactly
$\overline{U} \setminus U = \partial U$.  The general sub-additivity
inequality of Section~\ref{sec:valuation} still holds here
(\texttt{interiorMeasure\_add\_\allowbreak compl\_le}).  This
identification of slack with boundary mass also exhibits the modal
identification $\Bel = P(\Box\,\cdot)$ of the DS
literature~\cite{ruspini1986} arising from the logic itself.
Theorem~\ref{thm:finite-rep} is the
converse in the finite case (in the general case, the converse is the open
problem, see Section~\ref{sec:representation}).

To characterize weak excluded middle spaces, we give the 
(standard) definitions required: 

\begin{definition}[extremally disconnected]
A topological space $X$ is \emph{extremally disconnected} if the
closure of every open set is again open.
\end{definition}

\begin{definition}[Stone space, Stonean space]
The \emph{Stone space} of a Boolean algebra $B$ is the space of
ultrafilters on $B$ (equivalently, the $2$-valued homomorphisms
$B \to \{0,1\}$), topologized by the clopen basis $\{U : b \in U\}$
for $b \in B$. \emph{Stone duality}~\cite{johnstone1982} identifies
Boolean algebras with exactly the compact, Hausdorff, totally
disconnected spaces arising this way.  A \emph{Stonean space} is a compact Hausdorff space that
is also extremally disconnected. Equivalently, it is the Stone space
of a \emph{complete} Boolean algebra.
\end{definition}

\begin{remark}[\texttt{isDeMorgan\_opens\_iff\_\allowbreak
extremally\allowbreak Disconnected}]
Section~\ref{sec:valuation} noted that imposing weak excluded middle
on the whole frame collapses slack to pure double-negation
instability.  For $\Omega = \mathcal{O}(X)$, this global condition
holds exactly when $X$ is extremally disconnected.
\end{remark}

An immediate corollary gives the characterization:

\begin{corollary}[\texttt{isDeMorgan\_opens\_iff\_stonean}]
Given a compact Hausdorff space $X$, weak excluded middle on
$\mathcal{O}(X)$ holds exactly when $X$ is a Stonean space.
\end{corollary}

\subsection{The Computability Guard}
\label{sec:computability-guard}

The usefulness of the work done in this paper relies on the 
non-degeneracy of the spaces we study (i.e. ones where $\slack \equiv 0$
does not always hold). We must provide a witness for this non-degeneracy,
and we do so using computability and decidability concepts:

\begin{definition}[semidecidable proposition]
A proposition is \emph{semidecidable} if it is an open in the
Sierpi\'nski topology on \texttt{Prop}~\cite{escardo2004,rosolini1986}.
\end{definition}

\noindent
On the frame $\mathcal{O}(\mathbb{S}) =
\{\bot, \mathit{halts}, \top\}$ (containing three open sets),
the $\Sigma_1$ asymmetry is a theorem:

\begin{theorem}[\texttt{haltsOpen\_compl}]
\label{thm:sigma1-asymmetry}
$\neg\mathit{halts} = \bot$.
\end{theorem}

\noindent
Every nonempty open of $\mathbb{S}$ contains $\mathsf{True}$, so
$\bot$ is the only open disjoint from $\mathit{halts}$. Thus, no nonempty
open of $\mathbb{S}$ is evidence that the computation does not halt.

We now instantiate the \texttt{Measure.toValuationOpens} construction
of Section \ref{sec:interior-bridge} at $X = \mathbb{S}$,
allowing us to compute the slack:

\begin{theorem}[\texttt{slack\_eq\_boundary}]
\label{thm:slack-eq-boundary}
Let $\mu = p \cdot \delta_{\mathsf{True}} + (1-p) \cdot
\delta_{\mathsf{False}}$, and let $v$ be the resulting valuation.
Then
\[
  \slack(\mathit{halts}) \;=\; 1 - p \;=\;
  \mu\bigl(\partial\,\{q \mid q\}\bigr)
  \;=\; \mu\{\mathsf{False}\}.
\]
\end{theorem}

\noindent
Since $\mathit{halts}$ is already open, $v(\mathit{halts}) =
\mu(\mathit{halts}) = p$, which we can use to prove the above theorem. 
The measure of the slack (i.e. the measure of the outcome when the 
computation never halts) is the same as the measure of single point 
${\mathsf{False}}$ in a two-point space.
The DS ``ignorance'' of Section~\ref{sec:ds} is, in this model,
exactly this mass: the probability on the one outcome that no finite
observation can ever resolve. This probability is a fact about the model, 
not about anyone's ignorance.

\begin{definition}[halting problem]
For a code $c$ and input $n$, the \emph{halting problem} is the
predicate $H(c, n) = $ ``the machine with code $c$ halts on input
$n$,'' (\texttt{ComputablePred}). 
\end{definition}

\noindent
Slack cannot be eliminated by simply declaring a classical answer
for each code. The family of valuations
defined below attempts to do this:

\begin{definition}[sharp classical readout]
For a given $n$, the \emph{sharp classical readout} is the family
of slack-free valuations $\rho_n(c) \in \mathrm{Valuation}(\mathcal{O}(\mathbb{S}))$,
with $\rho_n(c)(\mathit{halts}) = 1$ if $H(c,n)$
holds, and $\rho_n(c)(\mathit{halts}) = 0$ otherwise.
\end{definition}

\noindent
However, this family is provably not computable:

\begin{theorem}[\texttt{sharpReadout\_not\_computable}]
For every $n$, the predicate $c \mapsto \bigl(\rho_n(c)(\mathit{halts}) = 1\bigr)$
is exactly $H(\cdot, n)$, hence not computable
(\texttt{ComputablePred.\allowbreak halting\_problem}).
\end{theorem}

\noindent
We can conclude that slack is not a defect
the constructive theory: any axiom strong enough to force classical 
collapse demands
of its models to compute the uncomputable. 
For the sum rule, this guard plays the role the
explicit \texttt{CoxModel} instance plays for the product rule
(Section~\ref{sec:cox}): each axiom set in the paper comes with 
non-trivial models satisfying it.
